\ifdefined\gkver\else\def\gkver{student}\fi
\documentclass[\gkver]{gaokaozhenti}
\newcommand{\ccnum}[1]{{\fontspec{Noto Sans CJK JP}#1}}
\begin{document}

\chapter{2013年全国新课标Ⅱ}
%% paperType: 理综物理部分

\begin{enumerate}
\item
%% number: 14
%% paperName: 2013年全国新课标Ⅱ第 14 题 6 分
%% typeId: 1
%% type: 单选题
%% chapter: 运动和力
%% point: 牛顿第二定律
%% method: 无
%% score: 6
%% degree: 650
%% duplicateId: 0
%% body:
一物块静止在粗糙的水平桌面上。从某时刻开始，物块受到一方向不变的水平拉力作用。假设物块与桌面间的最大静摩擦力等于滑动摩擦力。以 $a$ 表示物块的加速度大小，$F$ 表示水平拉力的大小。能正确描述 $F$ 与 $a$ 之间的关系的图像是\xzanswer{C}

\fourchoices[answer=C, ispicture=true, h=1.8cm]
{figs/14a.png}
{figs/14b.png}
{figs/14c.png}
{figs/14d.png}

%% answer: C
%% memo:
\memoanswer{
静摩擦力随外力而改变，当外力大于最大静摩擦力时，物体才产生加速度，可利用牛顿第二定律列方程求解。物块受到拉力和摩擦力作用，根据牛顿第二定律 $F-\mu mg=ma$，当 $F\leq F_{f\mathrm{max}}$ 时，$a=0$；当 $F>F_{f\mathrm{max}}$ 时，$a$ 与 $F$ 成一次函数关系，故选项 C 正确。
}

\item
%% number: 15
%% paperName: 2013年全国新课标Ⅱ第 15 题 6 分
%% typeId: 1
%% type: 单选题
%% chapter: 力物体的平衡
%% point: 摩擦力
%% method: 无
%% score: 6
%% degree: 650
%% duplicateId: 0
%% body:
如图，在固定斜面上的一物块受到一外力 $F$ 的作用，$F$ 平行于斜面向上。若要物块在斜面上保持静止，$F$ 的取值应有一定范围，已知其最大值和最小值分别为 $F_{1}$ 和 $F_{2}$（$F_{2}>0$）。由此可求出\xzanswer{C}

\onepicture[w=3.2cm, align=right]{figs/15.png}

\fourchoices[answer=C]
{物块的质量}
{斜面的倾角}
{物块与斜面间的最大静摩擦力}
{物块对斜面的正压力}

%% answer: C
%% memo:
\memoanswer{
本题原卷及材料中未提供详解，待补充。
}

\item
%% number: 16
%% paperName: 2013年全国新课标Ⅱ第 16 题 6 分
%% typeId: 1
%% type: 单选题
%% chapter: 电磁感应
%% point: 电磁感应中图象
%% method: 无
%% score: 6
%% degree: 650
%% duplicateId: 0
%% body:
如图，在光滑水平桌面上有一边长为 $L$、电阻为 $R$ 的正方形导线框；在导线框右侧有一宽度为 $d$（$d>L$）的条形匀强磁场区域，磁场的边界与导体框的一边平行，磁场方向竖直向下，导线框以某一初速度向右运动，$t=0$ 时导线框的右边恰与磁场的左边界重合，随后导线框进入并通过磁场区域。下列 $v-t$ 图像中，可能正确描述上述过程的是\xzanswer{D}

\onepicture[h=2.4cm, align=right]{figs/16.png}

\fourchoices[answer=D, ispicture=true, h=1.6cm]
{figs/16a.png}
{figs/16b.png}
{figs/16c.png}
{figs/16d.png}

%% answer: D
%% memo:
\memoanswer{
本题原卷及材料中未提供详解，待补充。
}

\item
%% number: 17
%% paperName: 2013年全国新课标Ⅱ第 17 题 7 分
%% typeId: 1
%% type: 单选题
%% chapter: 磁场
%% point: 带电粒子在磁场中的运动
%% method: 无
%% score: 7
%% degree: 650
%% duplicateId: 0
%% body:
空间有一圆柱形匀强磁场区域，该区域的横截面的半径为 $R$，磁场方向垂直横截面。一质量为 $m$、电荷量为 $q$（$q>0$）的粒子以速率 $v_{0}$ 沿横截面的某直径射入磁场，离开磁场时速度方向偏离入射方向 $60^\circ$。不计重力，该磁场的磁感应强度大小为\xzanswer{A}

\fourchoices[answer=A]
{$\frac{\sqrt{3}mv_{0}}{3qR}$}
{$\frac{mv_{0}}{qR}$}
{$\frac{\sqrt{3}mv_{0}}{qR}$}
{$\frac{3mv_{0}}{qR}$}

%% answer: A
%% memo:
\memoanswer{
本题原卷及材料中未提供详解，待补充。
}

\item
%% number: 18
%% paperName: 2013年全国新课标Ⅱ第 18 题 6 分
%% typeId: 1
%% type: 单选题
%% chapter: 电场
%% point: 电场叠加
%% method: 无
%% score: 6
%% degree: 650
%% duplicateId: 0
%% body:
如图，在光滑绝缘水平面上，三个带电小球 $a$、$b$ 和 $c$ 分别位于边长为 $l$ 的正三角形的三个顶点上；$a$、$b$ 带正电，电荷量均为 $q$，$c$ 带负电。整个系统置于方向水平的匀强电场中。已知静电力常量为 $k$。若三个小球均处于静止状态，则匀强电场场强的大小为\xzanswer{B}

\onepicture[w=2.8cm, align=right]{figs/18.png}

\fourchoices[answer=B]
{$\frac{\sqrt{3}kq}{3l^{2}}$}
{$\frac{\sqrt{3}kq}{l^{2}}$}
{$\frac{3kq}{l^{2}}$}
{$\frac{2\sqrt{3}kq}{l^{2}}$}

%% answer: B
%% memo:
\memoanswer{
本题原卷及材料中未提供详解，待补充。
}

\item
%% number: 19
%% paperName: 2013年全国新课标Ⅱ第 19 题 6 分
%% typeId: 2
%% type: 多选题
%% chapter: 电磁感应
%% point: 楞次定律
%% method: 无
%% score: 6
%% degree: 650
%% duplicateId: 0
%% body:
在物理学发展过程中，观测、实验、假说和逻辑推理等方法都起到了重要作用。下列叙述符合史实的是\xzanswer{ABD}

\fourchoices[answer=ABD]
{奥斯特在实验中观察到电流的磁效应，该效应解释了电和磁之间存在联系}
{安培根据通电螺线管的磁场和条形磁铁的磁场的相似性，提出了分子电流假说}
{法拉第在实验中观察到，在通有恒定电流的静止导线附近的固定导线圈中，会出现感应电流}
{楞次在分析了许多实验事实后提出，感应电流应具有这样的方向，即感应电流的磁场总要阻碍引起感应电流的磁通量的变化}

%% answer: ABD
%% memo:
\memoanswer{
本题原卷及材料中未提供详解，待补充。
}

\item
%% number: 20
%% paperName: 2013年全国新课标Ⅱ第 20 题 6 分
%% typeId: 2
%% type: 多选题
%% chapter: 曲线运动
%% point: 人造地球卫星
%% method: 无
%% score: 6
%% degree: 650
%% duplicateId: 0
%% body:
目前，在地球周围有许多人造地球卫星绕着它运转，其中一些卫星的轨道可近似为圆，且轨道半径逐渐变小。若卫星在轨道半径逐渐变小的过程中，只受到地球引力和稀薄气体阻力的作用，则下列判断正确的是\xzanswer{BD}

\fourchoices[answer=BD]
{卫星的动能逐渐减小}
{由于地球引力做正功，引力势能一定减小}
{由于气体阻力做负功，地球引力做正功，机械能保持不变}
{卫星克服气体阻力做的功小于引力势能的减小}

%% answer: BD
%% memo:
\memoanswer{
本题原卷及材料中未提供详解，待补充。
}

\item
%% number: 21
%% paperName: 2013年全国新课标Ⅱ第 21 题 6 分
%% typeId: 2
%% type: 多选题
%% chapter: 曲线运动
%% point: 圆周运动的应用
%% method: 无
%% score: 6
%% degree: 450
%% duplicateId: 0
%% body:
公路急转弯处通常是交通事故多发地带。如图，某公路急转弯处是一圆弧，当汽车行驶的速率为 $v_{c}$ 时，汽车恰好没有向公路内外两侧滑动的趋势。则在该弯道处\xzanswer{AC}

\onepicture[w=3.0cm, align=right]{figs/21.png}

\fourchoices[answer=AC]
{路面外侧高内侧低}
{车速只要低于 $v_{c}$，车辆便会向内侧滑动}
{车速虽然高于 $v_{c}$，但只要不超出某一最高限度，车辆便不会向外侧滑动}
{当路面结冰时，与未结冰时相比，$v_{c}$ 的值变小}

%% answer: AC
%% memo:
\memoanswer{
本题原卷及材料中未提供详解，待补充。
}

\item
%% number: 22
%% paperName: 2013年全国新课标Ⅱ第 22 题 8 分
%% typeId: 4
%% type: 实验题
%% chapter: 功和能
%% point: 机械能守恒定律
%% method: 无
%% score: 8
%% degree: 850
%% duplicateId: 0
%% body:
某同学利用下述装置对轻质弹簧的弹性势能进行探究：一轻质弹簧放置在光滑水平桌面上，弹簧左端固定，右端与一小球接触而不固连；弹簧处于原长时，小球恰好在桌面边缘，如图（a）所示。向左推小球，使弹簧压缩一段距离后由静止释放；小球离开桌面后落到水平地面。通过测量和计算，可求得弹簧被压缩后的弹性势能。

回答下列问题：
\begin{enumerate}
	\item
	本实验中可认为，弹簧被压缩后的弹性势能 $E_{p}$ 与小球抛出时的动能 $E_{k}$ 相等。已知重力加速度大小为 $g$。为求得 $E_{k}$，至少需要测量下列物理量中的\xzanswer{ABC}

	\fivechoices[answer=ABC]
	{小球的质量 $m$}
	{小球抛出点到落地点的水平距离 $s$}
	{桌面到地面的高度}
	{弹簧的压缩量 $\Delta x$}
	{弹簧原长 $l_{0}$}

	\item
	用所选取的测量量和已知量表示 $E_{k}$，得 $E_{k}=$\tkanswer[2.4cm]{$\frac{mgs^{2}}{4h}$}。

	\item
	图（b）中的直线是实验测量得到的 $s-\Delta x$ 图线。从理论上可推出，如果 $h$ 不变，$m$ 增加，$s-\Delta x$ 图线的斜率会\tkanswer[1.2cm]{减小}（填“增大”、“减小”或“不变”）；如果 $m$ 不变，$h$ 增加，$s-\Delta x$ 图线的斜率会\tkanswer[1.2cm]{增大}（填“增大”、“减小”或“不变”）。由图（b）中给出的直线关系和 $E_{k}$ 的表达式可知，$E_{p}$ 与 $\Delta x$ 的\tkanswer[1cm]{2}次方成正比。
\end{enumerate}

\twopicture[labelA=2013全国新课标Ⅱ22a, labelB=2013全国新课标Ⅱ22b, align=right, h=3cm]
{figs/22a.png}
{figs/22b.png}

%% answer: 
\jdanswer{
\begin{enumerate}
	\item
	ABC

	\item
	$\frac{mgs^{2}}{4h}$

	\item
	减小，增大，2

\end{enumerate}
}
%% memo:
\memoanswer{
本题原卷及材料中未提供详解，待补充。
}

\item
%% number: 23
%% paperName: 2013年全国新课标Ⅱ第 23 题 7 分
%% typeId: 4
%% type: 实验题
%% chapter: 电路
%% point: 电路实验
%% method: 无
%% score: 7
%% degree: 650
%% duplicateId: 0
%% body:
某同学用量程为 $1\UmA$、内阻为 $120\UO$ 的表头按图（a）所示电路改装成量程分别为 $1\UV$ 和 $1\UA$ 的多用电表。图中 $R_{1}$ 和 $R_{2}$ 为定值电阻，S 为开关。回答下列问题：
\begin{enumerate}
	\item
	根据图（a）所示的电路，在图（b）所示的实物图上连线。
	\item
	开关 S 闭合时，多用电表用于测量\tkanswer[1.6cm]{电流}（填“电流”、“电压”或“电阻”）；开关 S 断开时，多用电表用于测量\tkanswer[1.6cm]{电压}（填“电流”、“电压”或“电阻”）。
	\item
	表笔 A 应为\tkanswer[1.2cm]{黑}色（填“红”或“黑”）。
	\item
	定值电阻的阻值 $R_{1}=$\tkanswer[1.8cm]{$1.00\UO$}，$R_{2}=$\tkanswer[1.8cm]{$880\UO$}。（结果取 3 位有效数字）
\end{enumerate}

\onepicture[w=9cm, align=right]{figs/23.png}

%% answer: 
\jdanswer{
\begin{enumerate}
	\item
	如右图。

	\onepicture[w=7cm]{figs/23_an.png}

	\item
	电流，电压

	\item
	黑

	\item
	$1.00\UO$，$880\UO$

\end{enumerate}
}
%% memo:
\memoanswer{
本题原卷及材料中未提供详解，待补充。
}

\item
%% number: 24
%% paperName: 2013年全国新课标Ⅱ第 24 题 14 分
%% typeId: 6
%% type: 计算题
%% chapter: 电场
%% point: 带电粒子的圆周运动
%% method: 无
%% score: 14
%% degree: 650
%% duplicateId: 0
%% body:
如图，匀强电场中有一半径为 $r$ 的光滑绝缘圆轨道，轨道平面与电场方向平行。$a$、$b$ 为轨道直径的两端，该直径与电场方向平行。一电荷量为 $q$（$q>0$）的质点沿轨道内侧运动，经过 $a$ 点和 $b$ 点时对轨道压力的大小分别为 $N_{a}$ 和 $N_{b}$。不计重力，求电场强度的大小 $E$、质点经过 $a$ 点和 $b$ 点时的动能。

\onepicture[w=3.2cm, align=right]{figs/24.png}

%% answer: 
\jdanswer{
$E=\frac{1}{6q}(N_{b}-N_{a})$，$E_{ka}=\frac{r}{12}(N_{b}+5N_{a})$，$E_{kb}=\frac{r}{12}(5N_{b}+N_{a})$
}
%% memo:
\memoanswer{
质点所受电场力的大小为
\[f=qE\]
①

设质点质量为 $m$，经过 $a$ 点和 $b$ 点时的速度大小分别为 $v_{a}$ 和 $v_{b}$，由牛顿第二定律有
\[f+N_{a}=m\frac{v_{a}^{2}}{r}\]
②
\[N_{b}-f=m\frac{v_{b}^{2}}{r}\]
③

设质点经过 $a$ 点和 $b$ 点时的动能分别为 $E_{ka}$ 和 $E_{kb}$，有
\[E_{ka}=\frac{1}{2}mv_{a}^{2}\]
④
\[E_{kb}=\frac{1}{2}mv_{b}^{2}\]
⑤

根据动能定理有
\[E_{kb}-E_{ka}=2rf\]
⑥

联立①②③④⑤⑥式得
\[E=\frac{1}{6q}(N_{b}-N_{a})\]
⑦
\[E_{ka}=\frac{r}{12}(N_{b}+5N_{a})\]
⑧
\[E_{kb}=\frac{r}{12}(5N_{b}+N_{a})\]
⑨
}

\item
%% number: 25
%% paperName: 2013年全国新课标Ⅱ第 25 题 18 分
%% typeId: 6
%% type: 计算题
%% chapter: 运动和力
%% point: 牛顿定律的应用
%% method: 无
%% score: 18
%% degree: 450
%% duplicateId: 0
%% body:
一长木板在水平地面上运动，在 $t=0$ 时刻将一相对于地面静止的物块轻放到木板上，以后木板运动的速度-时间图像如图所示。已知物块与木板的质量相等，物块与木板间及木板与地面间均有摩擦，物块与木板间的最大静摩擦力等于滑动摩擦力，且物块始终在木板上。取重力加速度的大小 $g=10\Umsq$，求：
\begin{enumerate}
	\item
	物块与木板间、木板与地面间的动摩擦因数；
	\item
	从 $t=0$ 时刻到物块与木板均停止运动时，物块相对于木板的位移的大小。
\end{enumerate}

\onepicture[w=3.0cm, align=right]{figs/25.png}

%% answer: 
\jdanswer{
\begin{enumerate}
	\item
	$\mu_{1}=0.20$，$\mu_{2}=0.30$

	\item
	$s=1.125\Um$

\end{enumerate}
}
%% memo:
\memoanswer{
（1）从 $t=0$ 时开始，木板与物块之间的摩擦力使物块减速，此过程一直持续到物块和木板具有共同速度为止。

由图可知，在 $t_{1}=0.5\Us$ 时，物块和木板的速度相同，设 $t=0$ 到 $t=t_{1}$ 时间间隔内，物块和木板的加速度大小分别为 $a_{1}$ 和 $a_{2}$，则
\[a_{1}=\frac{v_{1}}{t_{1}}\]
①
\[a_{2}=\frac{v_{0}-v_{1}}{t_{1}}\]
②

式中 $v_{0}=5\Ums$，$v_{1}=1\Ums$ 分别为木板在 $t=0$、$t=t_{1}$ 时速度的大小。

设物块和木板为 $m$，物块和木板间、木板与地面间的动摩擦因数分别为 $\mu_{1}$ 和 $\mu_{2}$，由牛顿第二定律得
\[\mu_{1}mg=ma_{1}\]
③
\[(\mu_{1}+2\mu_{2})mg=ma_{2}\]
④

联立①②③④式得
\[\mu_{1}=0.20\]
⑤
\[\mu_{2}=0.30\]
⑥

（2）在 $t_{1}$ 时刻后，地面对木板的摩擦力阻碍木板运动，物块与木板之间的摩擦力改变方向。设物块与木板之间的摩擦力大小为 $f$，物块和木板的加速度大小分别为 $a_{1}'$ 和 $a_{2}'$，则由牛顿第二定律得
\[f=ma_{1}'\]
⑦
\[2\mu_{2}mg-f=ma_{2}'\]
⑧

假设 $f<\mu_{1}mg$，则 $a_{1}'=a_{2}'$；由⑤⑥⑦⑧式得 $f=\mu_{2}mg>\mu_{1}mg$，与假设矛盾，故
\[f=\mu_{1}mg\]
⑨

由⑦⑨式知，物块的加速度的大小 $a_{1}'$ 等于 $a_{1}$；物块的 $v-t$ 图象如图中点划线所示。

由运动学公式可推知，物块和木板相对于地面的运动距离分别为
\[s_{1}=2\times\frac{v_{1}^{2}}{2a_{1}}\]
⑩
\[s_{2}=\frac{v_{0}+v_{1}}{2}t_{1}+\frac{v_{1}^{2}}{2a_{2}'}\]
\ccnum{⑪}

物块相对于木板的位移的大小为
\[s=s_{2}-s_{1}\]
\ccnum{⑫}

联立①⑤⑥⑧⑨⑩\ccnum{⑪}\ccnum{⑫}式得 $s=1.125\Um$\ccnum{⑬}。
}

\item
%% number: 33
%% paperName: 2013年全国新课标Ⅱ第 33 题 10 分
%% typeId: 6
%% type: 计算题
%% chapter: 气体性质
%% point: 玻意耳定律
%% method: 无
%% score: 10
%% degree: 650
%% duplicateId: 0
%% body:
\begin{enumerate}
	\item
	（5 分）关于一定量的气体，下列说法正确的是\xzanswer{ABE}（填正确答案标号。选对 1 个得 2 分，选对 2 个得 4 分，选对 3 个得 5 分；每选错 1 个扣 3 分，最低得分为 0 分）。

	\fivechoices[answer=ABE]
	{气体的体积指的是该气体的分子所能到达的空间的体积，而不是该气体所有分子体积之和}
	{只要能减弱气体分子热运动的剧烈程度，气体的温度就可以降低}
	{在完全失重的情况下，气体对容器壁的压强为零}
	{气体从外界吸收热量，其内能一定增加}
	{气体在等压膨胀过程中温度一定升高}

	\item
	（10 分）如图，一上端开口、下端封闭的细长玻璃管竖直放置。玻璃管的下部封有长 $l_{1}=25.0\Ucm$ 的空气柱，中间有一段长 $l_{2}=25.0\Ucm$ 的水银柱，上部空气柱的长度 $l_{3}=40.0\Ucm$。已知大气压强为 $p_{0}=75.0\UcmHg$。现将一活塞（图中未画出）从玻璃管开口处缓慢往下推，使管下部空气柱长度变为 $l_{1}'=20.0\Ucm$。假设活塞下推过程中没有漏气，求活塞下推的距离。

	\onepicture[h=4.5cm, align=right]{figs/33.png}
\end{enumerate}

%% answer: 
\jdanswer{
\begin{enumerate}
	\item
	ABE

	\item
	$15.0\Ucm$

\end{enumerate}
}
%% memo:
\memoanswer{
以 $\UcmHg$ 为压强单位。在活塞下推时，玻璃管下部空气柱的压强为
\[p_{1}=p_{0}+l_{2}\]
①

设活塞下推后，下部空气柱的压强为 $p_{1}'$，由玻意耳定律得
\[p_{1}l_{1}=p_{2}l_{2}\]
②

如图，设活塞下推距离为 $\Delta l$，则此时玻璃管上部空气柱的长度为
\[l_{3}'=l_{3}+l_{1}-l_{1}'-\Delta l\]
③

设此时玻璃管上部空气柱的压强为 $p_{3}'$，则
\[p_{3}'-p_{1}'-l_{2}\]
④

由玻意耳定律得
\[p_{0}l_{3}=p_{3}'l_{3}'\]
⑤

由①至⑤式及题给数据解得 $\Delta l=15.0\Ucm$⑥。

\onepicture[w=3cm]{figs/33_an.png}
}

\item
%% number: 34
%% paperName: 2013年全国新课标Ⅱ第 34 题 10 分
%% typeId: 6
%% type: 计算题
%% chapter: 光学
%% point: 光的折射
%% method: 无
%% score: 10
%% degree: 650
%% duplicateId: 0
%% body:
\begin{enumerate}
	\item
	（5 分）如图，一轻弹簧一端固定，另一端连接一物块构成弹簧振子，该物块是由 $a$、$b$ 两个小物块粘在一起组成的。物块在光滑水平面上左右振动，振幅为 $A_{0}$，周期为 $T_{0}$。当物块向右通过平衡位置时，$a$、$b$ 之间的粘胶脱开；以后小物块 $a$ 振动的振幅和周期分别为 $A$ 和 $T$，则 $A$\tkanswer[1.2cm]{$<$}$A_{0}$（填“$>$”、“$<$”或“$=$”），$T$\tkanswer[1.2cm]{$<$}$T_{0}$（填“$>$”、“$<$”或“$=$”）。

	\onepicture[w=4.5cm, align=right]{figs/34a.png}

	\item
	（10 分）如图，三棱镜的横截面为直角三角形 $ABC$，$\angle A=30^\circ$，$\angle B=60^\circ$。一束平行于 $AC$ 边的光线自 $AB$ 边的 $P$ 点射入三棱镜，在 $AC$ 边发生反射后从 $BC$ 边的 $M$ 点射出，若光线在 $P$ 点的入射角和在 $M$ 点的折射角相等，
	\begin{enumerate}
		\item
		求三棱镜的折射率；
		\item
		在三棱镜的 $AC$ 边是否有光线透出，写出分析过程。（不考虑多次反射）
	\end{enumerate}

	\onepicture[w=4cm, align=right]{figs/34b.png}
\end{enumerate}

%% answer: 
\jdanswer{
\begin{enumerate}
	\item
	$<$，$<$

	\item
	$n=\sqrt{3}$，三棱镜的 $AC$ 边无光线透出。

\end{enumerate}
}
%% memo:
\memoanswer{
（i）光路图如图所示，图中 $N$ 点为光线在 $AC$ 边发生反射的入射点。设光线在 $P$ 点的入射角为 $i$、折射角为 $r$，在 $M$ 点的入射角为 $r'$、折射角依题意也为 $i$，有
\[i=60^\circ\]
①

由折射定律有
\[\sin i=n\sin r\]
②
\[n\sin r'=\sin i\]
③

由②③式得
\[r=r'\]
④

$OO'$ 为过 $M$ 点的法线，$\angle C$ 为直角，$OO'\parallel AC$，由几何关系有
\[\angle MNC=r'\]
⑤

由反射定律可知
\[\angle PNA=\angle MNC\]
⑥

联立④⑤⑥式得
\[\angle PNA=r\]
⑦

由几何关系得 $r=30^\circ$⑧。

联立①②⑧式得
\[n=\sqrt{3}\]
⑨

（ii）设在 $N$ 点的入射角为 $i''$，由几何关系得
\[i''=60^\circ\]
⑩

此三棱镜的全反射临界角满足
\[n\sin\theta_{c}=1\]
\ccnum{⑪}

由⑨⑩\ccnum{⑪}式得 $i''>\theta_{c}$\ccnum{⑫}。

此光线在 $N$ 点发生全反射，三棱镜的 $AC$ 边没有光线透出。

\onepicture[w=6cm]{figs/34_an.png}
}

\item
%% number: 35
%% paperName: 2013年全国新课标Ⅱ第 35 题 10 分
%% typeId: 6
%% type: 计算题
%% chapter: 运动和力
%% point: 动量守恒定律
%% method: 无
%% score: 10
%% degree: 450
%% duplicateId: 0
%% body:
\begin{enumerate}
	\item
	（5 分）关于原子核的结合能，下列说法正确的是\xzanswer{ABC}（填正确答案标号。选对 1 个得 2 分，选对 2 个得 4 分，选对 3 个得 5 分；每选错 1 个扣 3 分，最低得分为 0 分）。

	\fivechoices[answer=ABC]
	{原子核的结合能等于使其完全分解成自由核子所需的最小能量}
	{一重原子核衰变成 $\alpha$ 粒子和另一原子核，衰变产物的结合能之和一定大于原来重核的结合能}
	{铯原子核（$\ce{^{133}_{55}Cs}$）的结合能小于铅原子核（$\ce{^{208}_{82}Pb}$）的结合能}
	{比结合能越大，原子核越不稳定}
	{自由核子组成原子核时，其质量亏损所对应的能量大于该原子核的结合能}

	\item
	（10 分）如图，光滑水平直轨道上有三个质量均为 $m$ 的物块 $A$、$B$、$C$。$B$ 的左侧固定一轻弹簧（弹簧左侧的挡板质量不计）。设 $A$ 以速度 $v_{0}$ 朝 $B$ 运动，压缩弹簧；当 $A$、$B$ 速度相等时，$B$ 与 $C$ 恰好相碰并粘接在一起，然后继续运动。假设 $B$ 和 $C$ 碰撞过程时间极短。求从 $A$ 开始压缩弹簧直至与弹簧分离的过程中，
	\begin{enumerate}
		\item
		整个系统损失的机械能；
		\item
		弹簧被压缩到最短时的弹性势能。
	\end{enumerate}

	\onepicture[w=6cm, align=right]{figs/35.png}
\end{enumerate}

%% answer: 
\jdanswer{
\begin{enumerate}
	\item
	$\frac{1}{16}mv_{0}^{2}$

	\item
	$\frac{13}{48}mv_{0}^{2}$

\end{enumerate}
}
%% memo:
\memoanswer{
（i）从 $A$ 压缩弹簧到 $A$ 与 $B$ 具有相同速度 $v_{1}$ 时，对 $A$、$B$ 与弹簧组成的系统，由动量守恒定律得
\[mv_{0}=2mv_{1}\]
①

此时 $B$ 与 $C$ 发生完全非弹性碰撞，设碰撞后的瞬时速度为 $v_{2}$，损失的机械能为 $\Delta E$。对 $B$、$C$ 组成的系统，由动量守恒和能量守恒定律得
\[mv_{1}=2mv_{2}\]
②
\[\frac{1}{2}mv_{1}^{2}=\Delta E+\frac{1}{2}(2m)v_{2}^{2}\]
③

联立①②③式得
\[\Delta E=\frac{1}{16}mv_{0}^{2}\]
④

（ii）由②式可知 $v_{2}<v_{1}$，$A$ 将继续压缩弹簧，直至 $A$、$B$、$C$ 三者速度相同，设此时速度为 $v_{3}$，此时弹簧被压缩至最短，其弹性势能为 $E_{p}$。由动量守恒和能量守恒定律得
\[mv_{0}=3mv_{3}\]
⑤
\[\frac{1}{2}mv_{0}^{2}-\Delta E=\frac{1}{2}(3m)v_{3}^{2}+E_{p}\]
⑥

联立④⑤⑥式得
\[E_{p}=\frac{13}{48}mv_{0}^{2}\]
⑦
}

\end{enumerate}

\end{document}
